% Partie : counterfactuals
% Chapitre : minimal-change-semantics
% Section : transitivity

\documentclass[../../../include/open-logic-section]{subfiles}

\begin{document}

\olfileid[fr]{cnt}{min}{tra}

\olsection{Transitivité}

Pour le conditionnel matériel, la règle de chaînage est valide :
$!A \lif !B, !B \lif !C \Entails !A \lif !C$. Autrement dit, le
conditionnel matériel est transitif. En va-t-il de même des
conditionnels contrefactuels ? Considérons l'exemple suivant, dû à
Stalnaker :
\begin{quote}
  Si J.~Edgar Hoover était né en Russie, il aurait été communiste.

  Si J.~Edgar Hoover avait été communiste, il aurait été un traître.

  Donc, si J.~Edgar Hoover était né en Russie, il aurait été un
  traître.
\end{quote}
Si Hoover était né à la même date qu'en réalité, mais en Russie
plutôt qu'aux États-Unis, il y aurait grandi et serait devenu
communiste, supposons-le. La première prémisse est donc vraie. La
seconde prémisse, prise isolément, est également vraie. La conclusion
est pourtant fausse : selon toute vraisemblance, s'il était né en
Russie, Hoover serait devenu un communiste convaincu, et non un
traître à son pays. L'attribution intuitive de ces valeurs de vérité
s'accorde avec l'analyse de Stalnaker et Lewis. Parmi les mondes
possibles qui ne diffèrent du nôtre que par le lieu de naissance de
Hoover, le plus proche est celui où il devient un citoyen loyal de
l'URSS. Dans ce monde, le plus proche où l'antécédent de la première
prémisse et celui de la conclusion sont vrais, Hoover est un membre
loyal du parti communiste, et non un traître. Pour évaluer la seconde
prémisse, il faut au contraire considérer un autre monde : le plus
proche où Hoover est communiste, bien qu'il soit né aux États-Unis.
Il se serait rallié au communisme et serait donc, selon l'hypothèse de
l'exemple, un traître.\footnote{Pour apprécier la portée de
  l'exemple, il faut bien sûr accepter certaines hypothèses
  métaphysiques et politiques : par exemple, qu'il aurait été
  possible que Hoover naisse de parents russes, ou que les
  communistes vivant aux États-Unis dans les années 1950 aient été
  des traîtres à leur pays.}

\begin{prob}
  Trouvez un exemple intuitif convaincant où la transitivité des
  conditionnels contrefactuels échoue.
\end{prob}

\begin{ex}\ollabel{ex:trans-counterex}
  La sémantique des sphères invalide cette inférence : on a $p
  \cif q, q \cif r \Entails/ p \cif r$. Considérons le modèle
  $\mModel{M} = \tuple{W, O, V}$ où $W = \{w, w_1, w_2\}$,
  $O_w = \{\{w\}, \{w, w_1\}, \{w, w_1, w_2\}\}$,
  $V(p) = \{w_2\}$, $V(q) = \{w_1, w_2\}$ et
  $V(r) = \{w_1\}$. % Précision facultative conservée hors du texte traduit.
%
   Il existe une sphère admettant~$p$,
  $S = \{w, w_1, w_2\}$, et $p \lif q$ est vrai dans tous ses mondes ;
  donc $\mSat{M}{p \cif q}[w]$. Il existe aussi une sphère
  admettant~$q$, $S' = \{w, w_1\}$, et $q \lif r$ est vrai dans tous
  ses mondes ; donc $\mSat{M}{q \cif r}[w]$. En revanche, toute
  sphère admettant~$p$ contient~$w_2$, puisque c'est le seul monde
  où~$p$ est vrai. Or $\mSat/{M}{p \lif r}[w_2]$. Il existe
  bien une sphère admettant~$p$, mais aucune ne vérifie le
  conditionnel matériel en tous ses mondes. Le contrefactuel
  de la conclusion est donc faux au monde central.
\end{ex}

\begin{prob}
  Dessinez le système de sphères correspondant au contre-exemple de
  l'\olref[cnt][min][tra]{ex:trans-counterex}.
\end{prob}

\begin{prob}
  Dans l'\olref[cnt][min][tra]{ex:trans-counterex}, le monde~$w_2$ est
  celui où Hoover naît en Russie, devient communiste et n'est pas un
  traître ; $w_1$ est celui où il naît aux États-Unis, devient
  communiste et est un traître. Dans ce modèle, $w_1$ est plus proche
  de~$w$ que ne l'est~$w_2$. Est-ce nécessaire ? Pouvez-vous donner
  un contre-exemple sans supposer que la naissance de Hoover en
  Russie soit une possibilité plus éloignée que son adhésion au
  communisme ?
\end{prob}

\end{document}
